% =========================================================================
% CHAPTER 11 — CONCLUSION
% =========================================================================

\chapter{Conclusion}
\label{ch:conclusion}

\noindent
{\rmfamily\itshape
\begin{tabular}[t]{@{}l@{\hspace{0.3cm}}p{12.5cm}@{}}
\textls[50]{\textbf{Addressed:}} &
\textls[50]{Synthesis · Contribution · Future work}
\end{tabular}
}

\vspace{0.5cm}

This thesis began with an apparently simple question: how do you make a system understandable when it cannot be seen and to what extent machine learning-powered experiential representation can support the cognitive process of understanding?

The Atlantic Meridional Overturning Circulation served as the case, but the question was clearly never about the AMOC alone. It was about the wider class of systems --- invisible, dynamic, unfolding over timescales longer than a human life --- whose behaviour matters and whose structure resists the conventional ways we draw it. With this closing chapter, I tried to draw the arc together, statee what the work contributes, and mark the paths it opens.

% -------------------------------------------------------------------------
\section{Answering the Research Question}
\label{sec:concl-answer}

The three stages --- detect, translate, evaluate --- made different kinds of claims (Chapter~\ref{ch:rqs}) and therefore closed differently; that is part of the answer.

The \textbf{Detect} stage (RQ1) produced claims that could be checked against the data --- and they held: the vertical structure of the AMOC reduces to two modes explaining almost ninety-eight per cent of the variance, an autoencoder confirmed that the reduced space is sufficiently flat to serve as a coordinate system (H2), and the early-warning analysis returned a rigorous null result rather than a convenient signal (H3). These are the thesis's strongest findings: statements about a dataset.

The \textbf{Translate} stage (RQ2) asked only that the artefact exist and work. \emph{Living Atlantic} does: it runs and keeps its two conditions informationally equal through a shared data bundle and reconstruction (Chapter~\ref{ch:translate}).

The \textbf{Evaluate} stage (RQ3) was more cautious, asking how participants expressed understanding rather than whether they scored higher. Experiential presentation did not improve comprehension in every case; instead, it changed how participants understood the system and which aspects they could identify (Chapter~\ref{ch:discussion}). The two hypotheses therefore receive directional support rather than proof. H-A --- that conventional representations can leave structures inaccessible to non-specialists --- and H-B --- that experiential representation can make them more accessible --- find their clearest support in the recurrence scene, where no participant in the static condition identified the recurring structure, while several identified what was missing; support is weaker elsewhere, depending on whether a structure is temporal and invisible in a still figure or already visible in it.

The honest answer is not that experiential visualisation works, but that it works \emph{conditionally}: its usefulness depends on what needs to be understood, helping where a pattern is hard to see and offering less where it is already visible.

As Chapter~\ref{ch:introduction} anticipated, the main contribution is methodological, not oceanographic: a transferable framework, evidenced here through AMOC, for turning a complex, hard-to-observe system into a testable representation, with each stage held to its own standard of evidence --- verification, construction, and exploratory comparison.

Two further contributions sit alongside it. Empirically, the study surfaced a recurring misconception: more than half of participants read the height of the cumulative streamfunction curve as local velocity, a misreading that added richness did not reduce and that was rarer among more experienced participants --- unanticipated, and worth testing at scale. Materially, \emph{Living Atlantic} remains a reproducible design object built on informational equivalence, usable in future studies rather than only as a demonstration.

Together, these position the thesis as a bridge between scientific data and human understanding rather than a contribution to oceanography --- serving quality education and science communication (SDG~4) most directly, and, through its focus on the AMOC, climate action and ocean protection (SDG~13 and SDG~14).

% -------------------------------------------------------------------------
\section{Future Work}
\label{sec:concl-future}

The clearest next step follows from the findings of this study. The small convenience sample and differences in prior visualisation literacy limit the strength of the conclusions. A larger study could address these limitations by using a bigger and more balanced sample, while also considering participants' previous experience with data visualisations. This would make it possible to test more clearly whether experiential visualisation is more useful for some types of structures than others. The three scenes could therefore be treated as separate tests of the same broader question --- when does experiential presentation actually improve understanding?

The misconception identified in the study also suggests a clear design improvement. Because the height of the streamfunction curve can easily be interpreted as speed, a future version should make its cumulative nature more visible. One option would be an animation that builds the integral as the cursor moves down the water column, while showing the local slope alongside it as the quantity related to direction and strength. This would make the existing mathematics easier to understand and directly address the most common misconception found in the study.

More broadly, the framework could be applied beyond the AMOC. Many important systems are difficult to observe directly, so testing the same approach with other complex systems could show whether it can also help make invisible structures easier to understand.

% -------------------------------------------------------------------------
\section{Closing}
\label{sec:concl-closing}

The thesis set out to make an invisible system intelligible, and to ask whether changing how we experience data changes how we understand it. It leaves that question not answered once and for all, but answered with precision: for some structures, for some viewers, and at a perceptual cost that must be designed around. The invisible did not become fully visible. But it became, in places and for reasons the study can name, a little more intelligible --- and it showed exactly where the next attempt should aim.